\section{Refined crosscap invariants}\label{sec:setup}

\subsection{Variables}\label{sec:variables}
\begin{center}
\small
\begin{tabular}{>{$}l<{$}p{9.5cm}}
\toprule
\text{symbol} & meaning \\
\midrule
q,\ t & variables of refined Chern--Simons theory (as in \cite{KN17}); half-integer powers $q^{1/2}, t^{1/2}$ are written as such \\
\lambda = t^{N - 1/2} & Kauffman variable ($\lambda$ of \cite{NRZ13}, $\nu$ of \cite{Marino09}); $N$ is the rank of $D_N$ and $t^N = \lambda t^{1/2}$. At $t = q$ it is the $\lambda = q_{\mathrm{NRZ}}^{2N-1}$ of \cite{NRZ13} ($q_{\mathrm{NRZ}} = q^{1/2}$); at large $N$ it is treated as a formal variable \\
\tlam = \lambda\, \qr & shifted Kauffman variable ($-\ba\bt$ in bold variables) \\
a = \lambda^2 = t^{N_A},\ N_A = 2N - 1 & K\"ahler variable of the cover $\U(N_A)$ ($a = t^{N_A}$ of \cite{KN17}) \\
(\bq, \bt, \ba) = (t^{1/2}, -\qr, \lambda) & bold variables; $\bq, \bt$ as in \cite{KN17} ($q^{1/2} = -\bq\bt$, $t^{1/2} = \bq$), $\ba = \lambda$ as in \cite{AgSc12}. The $\ba$ of the cover in \cite{KN17} is defined by $a = -\ba_{\mathrm{KN}}^2\bt$, so $\ba_{\mathrm{KN}} = \ba\,(t/q)^{1/4}$ \\
\ql = (qt)^{1/2},\ \qr = (q/t)^{1/2} & $q_\ell, q_r$ of \cite{KN17} ($q = \ql\qr$, $t = \ql/\qr$); $w := \qr$ is the notation of \cite{AM22} \\
z_q = q^{1/2} - q^{-1/2},\ z_t = t^{1/2} - t^{-1/2} & $z_q^+ = q^{1/2} + q^{-1/2}$ \\
G = z_q z_t = \ql + \ql^{-1} - \qr - \qr^{-1} & variable of the $G$-expansion; $(\bq\bt - \bq^{-1}\bt^{-1})(\bq - \bq^{-1}) = -G$ in \cite{KN17} \\
\sigma\colon q \mapsto t^{-1},\ t \mapsto q^{-1} & $\ql \mapsto \ql^{-1}$ with $\qr$ and $\lambda$ fixed; $G$ is invariant \\
\text{mirror} & $\lambda \mapsto \lambda^{-1}$, $q \mapsto q^{-1}$, $t \mapsto t^{-1}$ (the mirror image of the knot) \\
\delta & the Seifert genus of the knot: $\delta = m$ for $T(2,2m+1)$, $(r-1)(s-1)/2$ for $T(r,s)$; for thin knots it equals the $\delta$-degree $\deg_{\ba} + \tfrac12\deg_{\bq} - \deg_{\bt}$ of the superpolynomial \\
\bottomrule
\end{tabular}
\end{center}
In the unrefined limit $t = q$ the variables of \cite{Marino09} are $t_M = q^{1/2}$ and $z = z_q$, and the $q_{\mathrm{NRZ}}$ of the Kauffman polynomial $F_R(K; \lambda, q_{\mathrm{NRZ}})$ of \cite{NRZ13} is also $q^{1/2}$. Young diagrams are written as $R, \mu, \nu$ (never $\lambda$); in displayed formulas we draw them, e.g.\ $g_{\ytwo}$, and in the text we write $[2]$, $(2)$ for the same diagram. Composite representations of $\U(N_A)$ are written $(\alpha;\bar\beta)$; the adjoint is $(1;\bar1)$.

\subsection{Refined \texorpdfstring{$\SO(2N)$}{SO(2N)} invariants}
In the $\SO(2N)$ refined Chern--Simons theory of \cite{AgSc12}, the Wilson loop in a representation $R$ is given by a Macdonald polynomial of type $D_N$.
% [S1]
\begin{result}\label{cr:dim}
The refined unknot $\dim_R$ of \cite{AgSc12} equals the principal specialization $x = (1, t, \dots, t^{N-1})$ of the Koornwinder polynomial \cite{Koornwinder92} $K_R(x; q, t; 1, -1, q^{1/2}, -q^{1/2})$ for six representations with $|R| \le 3$ (checked against the evaluation formulas of \cite{vD96,Sahi99} and against polynomials constructed by the $q$-difference operator for $N = 2, 3$).
$\dim_R$ is not a Laurent polynomial in $t^N$: for instance, for the representation $[2]$ it has the denominator $(1 - qt^{N-1})$ and for $[1,1]$ the denominator $(1 - qt^{2N-3})$; both cancel at $t = q$.
\end{result}

% [S2]
\begin{remark}[failure of the naive refinement]
If the construction of \cite{Marino09} is refined as it stands (bosonic Cauchy kernel, the factor $\tfrac12$, the sum over odd $d$), the poles in $U = t^N$ fail to cancel already for the two-box unknot. The composite $([1],[1])$ contributes only in the $p_1^2$ direction, whereas the poles of the $D_N$ representations $[2]$ and $[1,1]$ sit at different places ($t^N = t/q$ and $t^{2N} = t^3/q$). None of five choices of weights and three choices of basis removes the mismatch.
\end{remark}

\subsection{Dual Cauchy kernel and generating function}\label{sec:kernel}
% [S3]
In the dual Cauchy identity for Koornwinder polynomials \cite{Mimachi01},
\begin{equation}\label{eq:mimachi}
\sum_{\mu \subseteq m^n} (-1)^{|\mu|} K_{m^n - \mu}(x; q, t)\, K_{\mu'}(y; t, q) = \prod_{i, j} x_i^{-1}(1 - x_i y_j)(1 - x_i y_j^{-1}),
\end{equation}
the holonomy side carries the $D_n$-type parameters $(\pm1, \pm q^{1/2})$ of \cite{AgSc12} and the probe side the parameters $(\pm 1, \pm t'^{1/2})$ for $(q', t') = (t, q)$. We checked the identity for $(n, m) = (1,1), (2,1), (1,2), (2,2), (2,3), (3,2)$.
We read this identity as the Ooguri--Vafa kernel of the refined orientifold and set, for probe variables $Y = (y_1, y_2, \dots)$,
\begin{equation}\label{eq:Z}
\tilde Z(Y) = \Bigl\langle \prod_{\alpha, i} (1 - y_\alpha x_i)(1 - y_\alpha x_i^{-1}) \Bigr\rangle_{\SO(2N)},\qquad
\log \tilde Z = \sum_{d \ge 1} \frac1d \sum_R f_R(q^d, t^d, \lambda^d)\, s_R(Y^d).
\end{equation}
Here $\langle F(x)\rangle_{\SO(2N)}$ is the linear map that expands a class function $F$ of the holonomy eigenvalues $x = (x_1, \dots, x_N)$ in Koornwinder polynomials $K_\mu(x; q, t; 1, -1, q^{1/2}, -q^{1/2})$ and replaces each term by the refined Chern--Simons knot invariant $W_\mu(K)$ (normalized as in \S\ref{sec:uniform}): $F = \sum_\mu c_\mu K_\mu(x)$ gives $\langle F\rangle := \sum_\mu c_\mu W_\mu(K)$. In the unrefined theory this is the path-integral expectation value of the Wilson loop, in which $x$ is integrated out. The sum in \eqref{eq:Z} runs over all $d$ (``all-$d$''), not only over odd $d$ as in \cite{Marino09}.

Since $(1 - x_iy_j^{-1}) = -x_iy_j^{-1}(1 - y_jx_i^{-1})$, the right-hand side of \eqref{eq:mimachi} with $n = N$ equals $(-1)^{Nm}(\prod_j y_j)^{-N}\prod_{i,j}(1 - y_jx_i)(1 - y_jx_i^{-1})$. Reading the identity with $m$ probes on the $y$ side and $N$ holonomies on the $x$ side, the kernel of \eqref{eq:Z} is therefore
\begin{equation}\label{eq:Zexpand}
\begin{aligned}
\tilde Z(Y) &= \sum_{\mu \subseteq m^N} (-1)^{|\mu| + Nm}\, W_\mu(K)\, \hat K_{\mu'}(Y; t, q; T = t^N),\\
\hat K_{\mu'}(Y; t, q; T) &:= \Bigl(\prod_\alpha y_\alpha\Bigr)^N K_{N^m - \mu'}(Y; t, q; 1, -1, q^{1/2}, -q^{1/2}).
\end{aligned}
\end{equation}
The expansion runs over Koornwinder polynomials of type $(t, q)$ on the probe side; the complement in the rectangle $N^m$ multiplied by $(\prod y)^N$ is a virtual Koornwinder polynomial \cite{Rains05,RW21}, which depends on $N$ only through $T = t^N = U$. The kernel in \eqref{eq:Z} has no prefactor, so setting $y_{m+1} = 0$ reduces the $(m+1)$-probe kernel to the $m$-probe one; hence $\log\tilde Z$ is a well-defined symmetric function of $Y$ and $f_R$ does not depend on the number of probes once $m \ge \ell(R)$. The invariant $f_R$ is defined by expanding $\log\tilde Z$ in Schur functions $s_R(Y)$ and applying M\"obius inversion in $d$. In the actual computation (Appendix \ref{app:code}) the kernel $\prod_{\alpha,i}(1 - y_\alpha x_i)(1 - y_\alpha x_i^{-1}) = \exp(-\sum_d p_d(Y)S_d(x)/d)$, $S_d(x) = \sum_i(x_i^d + x_i^{-d})$, is expanded in $K_\mu(x)$ by evaluating at the points $x = t^\rho q^\nu$, where $S_d(t^\rho q^\nu)$ is a rational function of $U$; this is why $U$ can be treated as a formal variable, independently of the closed form of \S\ref{sec:virtual}. The sign $(-1)^{|\mu|}$ is the hallmark of a fermionic kernel. For the unknot $W_\mu = P_\mu(t^\rho)$ and \eqref{eq:Zexpand} reduces to the dual Cauchy identity evaluated at $x = t^\rho$, which is a product (\S\ref{sec:unknot}).

For the cover $\U(N_A)$ we define $f^{\cov}_R$ with the same kernel $\prod_i(1 - y_\alpha x_i)(1 - y_\alpha x_i^{-1})$. Since both $x$ and $x^{-1}$ enter, the expansion basis consists of Macdonald characters of composite representations $(\alpha; \bar\beta)$, and $\langle P_\alpha(x) P_\beta(x^{-1})\rangle$ is evaluated by decomposing into composite values and contraction terms (\S\ref{sec:generalN}).
The part of the kernel with $|Y| \le 2$ is expressed through ordinary representations of $\U(N_A)$ and the adjoint composite $(1;\bar1)$; at degree three the composites $(2;\bar1)$ and $(1,1;\bar1)$ appear as well.

\begin{definition}[refined crosscap invariant]\label{def:crosscap}
In the all-$d$ definition \eqref{eq:Z} we set $g_R = f_R - \tfrac12 f^{\cov}_R$.
The hatted invariant is $\hat g = M(q^{-1})^{-1} g$, where $M_{RS} = \sum_\mu \chi_R(\mu)\chi_S(\mu) z_\mu^{-1}\prod_i (t_M^{\mu_i} - t_M^{-\mu_i})/(t_M - t_M^{-1})$ with $t_M = q^{1/2}$ is the matrix $M$ of \cite{Marino09}, which equals $M(x)$ of \cite{KN17} at $x = q^{-1}$.
\end{definition}

% [S8]
\begin{remark}
With $M(t)$ or $M(t^{-1})$ in place of $M(q^{-1})$ the Hopf link invariants are not Laurent polynomials; $M(q)$ exchanges $R \leftrightarrow R^T$.
The two formulations of \cite{KN17} (the $q$-side with the Macdonald Cauchy kernel and $M(t)$, the $t$-side with the dual Cauchy kernel and $M(q^{-1})$) give the same hatted invariants with two boxes (\S\ref{sec:generalN}).
\end{remark}

\subsection{Computational scheme}\label{sec:procedure}
The invariant $g_R$ of Definition \ref{def:crosscap} is computed in three steps. $f_R$ and $f^{\cov}_R$ are obtained separately, order by order in the number of boxes, and $g_R$ is their difference.
\begin{center}
\small
\begin{tabular}{lp{6.6cm}p{5.4cm}}
\toprule
quantity & input & computation \\
\midrule
$f_R$ & $\SO(2N)$ refined Chern--Simons invariants $W^{\SO}_\mu$, $|\mu| \le |R|$, type $D_N$ & expand $\log$ of the dual Cauchy kernel \eqref{eq:Z} in Schur functions and apply M\"obius inversion \\
$f^{\cov}_R$ & A-type invariants of the cover $\U(N_A)$: ordinary representations $W_\alpha$, the composite $(1;\bar 1)$ (adjoint), and for three boxes $(2;\bar 1), (1,1;\bar 1)$ with contraction terms & same kernel, expanded in composite Macdonald characters \\
$g_R$ & the two above & $f_R - \tfrac12 f^{\cov}_R$; hatted with $M(q^{-1})^{-1}$ \\
\bottomrule
\end{tabular}
\end{center}
$f_R$ closes on the $\SO$ side, and $f^{\cov}_R$ contains no $\SO$ data. That the cover is given by A-type refined invariants is a hypothesis---the refined version of the construction of \cite{Marino09}; for the fundamental representation it is stated as Hypothesis~\ref{hyp:hg}. It is supported by the fact that the cancellation of $N$-dependent poles after the subtraction fixes the framing monomial uniquely (Result~\ref{cr:uniform} and Result~\ref{cr:monomial}) and by the integrality of $g_R$ (\S\ref{sec:generalN}).
Writing the cover with two sets of probes, brane ($Y$) and image ($Y'$), $f^{\cov}$ splits into a brane--brane part (the invariant of \cite{KN17} itself) and a brane--image part (the connected part of the composites), see \S\ref{sec:cover}.

\subsection{Uniform normalization}\label{sec:uniform}
The input $\U(n)$ refined invariants, normalized by conditions (N1) and (N2) of \cite{KN17}, change by powers of $(q/t)$ under determinant shifts and charge conjugation. With $\delta$ the Seifert genus (\S\ref{sec:variables}) we set
\begin{equation}\label{eq:uniform}
W_R = (t/q)^{\delta |R|/2}\, \Rred_R\, P_R(t^\rho),
\end{equation}
where $\Rred_R$ is the reduced invariant (the refined invariant divided by that of the unknot; in \S\ref{sec:generalN} it is computed from DAHA-Jones polynomials) and $P_R(t^\rho)$ the Macdonald dimension.

% [S5][S6]
\begin{result}\label{cr:uniform}
For the normalization \eqref{eq:uniform} the following hold.
\begin{enumerate}
\item $W$ is exactly invariant under the determinant shift $R \to R + 1^n$ and under charge conjugation. Checked for 16 pairs for the trefoil ($\U(2)$--$\U(5)$), 13 pairs each for $T(2,5)$ and $T(2,7)$, and seven double twist knots.
\item At fixed $N$ ($N = 2, 3$, $|R| \le 3$), fixing the framing monomial by requiring the cancellation of $N$-dependent poles and the value $1$ at $t = q$ gives the monomial $1$ on both the $\SO$ side and the cover side ($T(2,3), T(2,5), T(2,7)$). Within the family $(t/q)^{\beta/2}$ the exponent $\beta$ is unique.
\item The DAHA-Jones polynomials (vertical framing) agree up to powers of the framing eigenvalue (\S\ref{sec:generalN}).
\item At $t = q$ this is the standard framing of \cite{Marino09}.
\end{enumerate}
\end{result}

% [S7]
\begin{hypothesis}\label{hyp:framing}
The uniform normalization \eqref{eq:uniform} is the normalization at framing zero.
\end{hypothesis}
We use this normalization on both sides throughout.

\subsection{The unknot}\label{sec:unknot}
% [H1]
\begin{proposition}\label{prop:unknot}
For the unknot, $f_{10} = (1 - t^N)(t^N + t)/(t^N(t-1))$ and $f_{a0} = 0$ for $a \ge 2$ (the first index counts probes on the knot, the second on a second component, cf.\ \S\ref{sec:hopf}). In terms of $\lambda$,
\[
f_{10} = -1 - \frac{\lambda - \lambda^{-1}}{t^{1/2} - t^{-1/2}},
\]
the sum of the unknot of \cite{Marino09}, $(\lambda - \lambda^{-1})/z$, and the crosscap term $-1$. In our conventions $g_{\yone}$ is minus the mirror of the crosscap invariant of \cite{Marino09} (Result \ref{cr:marino}), so $-1$ corresponds to the value $g_{\yone} = 1$ of \cite{Marino09} for the unknot. The product formula is $\prod_j (y_j^{\pm}; t)_\infty / (U y_j^{\pm}; t)_\infty$; neither the factor $\tfrac12$ nor the restriction to odd $d$ is needed.
\end{proposition}
\begin{proof}
The expectation value of the dual Cauchy kernel is $\prod_j\prod_{k=1}^{N} t^{-(N-k)}(1 - t^{N-k}y_j^{\pm})$; take the logarithm.
\end{proof}
